% Cortes mecánicos íntegros: parte_iv.tex:55--106,337--408.
% Residencia material del ensamblaje sucesor de 102 capítulos.
% Residencia de realización: capítulo 76.
\section{Suspension spatiale}

Soit
\[
 \Gamma(\theta)=
 \bigl(r(\theta)\cos\theta,r(\theta)\sin\theta,z(\theta)\bigr).
\]
Sa courbure et sa torsion se calculent au moyen de
\[
 \kappa_{\mathrm{F}}=\frac{\|\Gamma'\times\Gamma''\|}{\|\Gamma'\|^3},
 \qquad
 \mathcal T_{\mathrm{F}}=
 \frac{\det(\Gamma',\Gamma'',\Gamma''')}
 {\|\Gamma'\times\Gamma''\|^2}.
\]
Une hélice générale au sens de Lancret satisfait
\[
 \frac{\mathcal T_{\mathrm{F}}}{\kappa_{\mathrm{F}}}=\text{constante}.
\]
Une suspension radiale arbitraire ne satisfait pas nécessairement cette condition. Le nom \emph{suspension hélicoïdale radiale} conserve la structure d'enroulement sans affirmer une propriété différentielle non démontrée.

\section{Trichromie géométrique}

Les générateurs
\[
 J_{+1}^2=-I,\qquad J_0^2=0,\qquad J_{-1}^2=I
\]
produisent des évolutions elliptique, parabolique et hyperbolique. Une histoire peut les alterner et publier une courbe dont la courbure de Frenet varie continûment. La \emph{trichromie} décrit la séquence des régimes internes ; la géométrie de Frenet décrit la sortie. Le lecteur radial conserve les deux sans les confondre.

\begin{figure}[htbp]
\centering
\begin{tikzpicture}[>=Latex,node distance=10mm and 16mm,
 n/.style={draw=HMTNavy,fill=white,text width=34mm,minimum height=10mm,align=center,font=\small\sffamily}]
\node[n] (int) {historia TRIT\\\(+1,0,-1,\ldots\)};
\node[n,right=of int] (rad) {lecteur radial\\\(p,\Lambda,C,\mu\)};
\node[n,right=of rad] (curve) {courbe publiée\\\(r(\theta),z(\theta)\)};
\node[n,below=of curve] (frenet) {géométrie de Frenet\\\(\kappa_{\mathrm{F}},\mathcal T_{\mathrm{F}}\)};
\draw[->,HMTBlue,thick] (int)--(rad);
\draw[->,HMTBlue,thick] (rad)--(curve);
\draw[->,HMTBlue,thick] (curve)--(frenet);
\draw[->,HMTRed,dashed,bend left=28] (frenet.west) to node[below,font=\scriptsize\sffamily]{ne sélectionne pas rétrospectivement} (int.south);
\end{tikzpicture}
\caption{Séparation causale entre régime interne et courbure visible.}
\label{espiral:fig:trit-frenet}
\end{figure}

\chapter{Structure dynamique de l'espace et conséquences}

\section{Lamination de feuillets hélicoïdaux}

La synthèse géométrique est une lamination dont les éléments sont des histoires relevées. Chaque feuillet porte :
\begin{enumerate}
\item une phase angulaire ;
\item une coordonnée radiale interne ;
\item une profondeur de mémoire ;
\item une séquence de feuillets APP ;
\item un régime TRIT local ;
\item une holonomie TPK et son bord.
\end{enumerate}
Les sections transversales sont des cercles ou des courbes de niveau ; les familles de sections forment des cylindres ; la variation radiale produit des suspensions ; l'involution produit des feuillets conjugués. La géométrie réelle est l'image de cette lamination sous une publication qui peut contracter plusieurs feuillets.

\section{Classification cinématique par incréments angulaire et radial}

Une droite radiale apparaît lorsque \(\omega=0\) et \(\beta\neq0\). Un cercle apparaît lorsque \(\omega\neq0\) et \(\beta=0\). La spirale apparaît lorsque les deux sont non nuls. Ces trois sorties sont des configurations d'un même mécanisme de transport, non des objets sans relation. Le support APP est bidimensionnel dès le départ ; le parcours unidimensionnel est une section particulière de ce support.

\section{Quotient de Möbius par retour avec inversion d'orientation}

Lorsque l'involution de voie est identifiée dans la couture après un tour, le quotient peut acquérir un ruban de Möbius. L'affirmation exige une action libre et l'opérateur d'entrelacement qui transporte l'orientation. Le ruban ne se déduit pas de la simple présence d'un signe \(-1\) ; il se déduit de la couture
\[
 (x,0)\sim(C_{\mathrm{rev}}x,1)
\]
dans le domaine déclaré. Dans l'atlas matériel, l'involution dodécaphasique et l'inversion cellulaire restent des opérations distinctes.

\Needspace{34\baselineskip}
\section{Classification de 3+2 points d'équilibre au moyen du hessien}

L'intuition des trois points colinéaires et des deux points triangulaires de Lagrange peut être formulée comme une preuve ultérieure. Étant donné un potentiel effectif \(V_{\mathrm{eff}}\) produit par la réalisation HMT--MD, les points critiques satisfont
\[
 \nabla V_{\mathrm{eff}}=0.
\]
Leur classification exige le hessien et l'indice de Morse. Ce n'est que si l'application dimensionnelle produit exactement trois points critiques colinéaires et deux triangulaires qu'elle peut être reliée à la partition \(3+2\). La correspondance n'est pas utilisée comme entrée de la taxonomie spirale.

\section{Hiérarchie causale entre courbes, champs et sections matérielles}

La correspondance spirale ouvre trois sorties coordonnées :
\[
 \text{registre APP--TRIT--TPK}
 \longrightarrow
 \begin{cases}
 \text{courbe archimédienne},\\
 \text{mode électromagnétique},\\
 \text{signature de route matérielle}.
 \end{cases}
\]
La première est fermée algébriquement dans ce travail. La seconde possède une réalisation exacte pour les modes circulaires et une condition de compatibilité spectrale. La troisième dispose de l'électron central et de l'atlas, mais exige de vérifier le carré de constance sur chaque multisection avant d'attribuer des taxonomies de saveur.
